\documentclass[11pt,a4paper]{article}

\usepackage[utf8]{inputenc}
\usepackage[T1]{fontenc}
\usepackage[margin=1in]{geometry}
\usepackage[expansion=false]{microtype}
\usepackage{titlesec}
\usepackage{enumitem}
\usepackage{xcolor}
\usepackage{hyperref}
\usepackage{abstract}

\definecolor{accent}{HTML}{1F4E5C}
\hypersetup{colorlinks=true, linkcolor=accent, citecolor=accent, urlcolor=accent}

\titleformat{\section}{\large\bfseries\color{accent}}{\thesection}{0.6em}{}
\titleformat{\subsection}{\normalsize\bfseries\color{accent}}{\thesubsection}{0.6em}{}
\titlespacing*{\section}{0pt}{1.4ex plus 1ex minus .2ex}{0.8ex plus .2ex}

\setlist[itemize]{leftmargin=1.4em, itemsep=2pt, topsep=3pt}
\renewcommand{\abstractnamefont}{\normalfont\bfseries\color{accent}}

\title{\vspace{-2.2em}\bfseries\color{accent} Predicting Per-Screen Time-on-Task from User Interface Screens}
\author{\normalsize Master's Thesis Proposal}
\date{\normalsize \today}

\begin{document}
\maketitle
\vspace{-2.2em}

\begin{abstract}
\noindent
Time-on-Task --- the time a user spends to accomplish a goal --- is the principal behavioural measure of the \emph{efficiency} dimension of usability defined in ISO~9241-11. Estimating it today still requires either live user testing or hand-built cognitive models, both costly and neither scaling to the millions of interfaces deployed in practice. This thesis investigates whether per-screen Time-on-Task can be predicted directly from the visual and structural properties of an interface screen, learned from large-scale real-world interaction traces and validated in two stages: against held-out interface domains never seen in training, and --- exactly once, zero-shot --- against human times measured in an independent study. The work positions itself as a data-driven successor to the classical predictive tradition of Keystroke-Level and GOMS modelling, removing those models' reliance on manual operator coding and their assumption of expert, error-free behaviour. Beyond a predictive model, the central scientific contribution is a validated methodology for recovering per-screen timing labels from interaction logs --- with confounds such as page-load latency identified and reported rather than absorbed --- and an honest characterisation of how much of human task time is recoverable from a screen alone.
\end{abstract}

\section{Motivation and Problem Statement}
ISO~9241-11 frames usability as effectiveness, efficiency, and satisfaction in a context of use, and in its 2018 revision defines efficiency explicitly in terms of the \emph{resources --- including time --- expended to achieve a goal}~\cite{iso9241}. Of the three dimensions, efficiency is the one most directly captured by a single objective measure: Time-on-Task (ToT). It is among the oldest and most widely reported usability metrics, collected in industrial usability testing since the 1980s and codified through the MUSiC project into the modern standard.

The difficulty is that obtaining ToT remains expensive. It is measured either empirically, by recruiting participants and timing them, or analytically, by constructing a predictive cognitive model of the interaction by hand. Neither approach scales to the volume and churn of contemporary digital interfaces, where designers would benefit from an early, automatic estimate of how much time a given screen will cost its users --- well before any user is recruited.

This thesis asks whether that estimate can be produced from the screen itself. A screen's layout, density, and structure plausibly carry much of the information that determines how long a user will dwell on it before acting. If so, a model that reads a screen and returns an expected dwell time would serve as a cheap, scalable instrument for usability screening, complementing rather than replacing empirical testing.

\section{Background and Related Work}

\paragraph{The predictive tradition in HCI.} The idea that interaction time can be predicted from an interface specification, without user testing, originates with the Keystroke-Level Model (KLM) of Card, Moran, and Newell~\cite{klm1980,psychologyhci}. KLM decomposes a task into primitive operators --- keystrokes, pointing, homing, mental preparation, and system response --- and sums their empirically calibrated durations to predict expert, error-free completion time, reportedly to within roughly 20\%. As the simplest member of the GOMS family, it remains a standard usability-engineering tool. Its limitations are equally well known and directly motivate this work: it requires an analyst to manually specify the operator sequence for every task, it assumes a skilled user executing a known method, and it does not account for the visual layout or readability of the interface~\cite{klmreview}.

\paragraph{Learned models of interface performance.} More recent work replaces hand-coded operators with learned predictors. Models in the gradient-descent UI-optimization line take a structured representation of a user interface and a task sequence and predict per-task completion time, using recurrent architectures to capture the learning effect across a session~\cite{deepui}. These models demonstrate that task time is learnable, but they operate on constrained interfaces (e.g.\ menus, synthetic mobile layouts) and require structured UI descriptions rather than the rendered screen.

\paragraph{Visual search and attention on GUIs.} A parallel line predicts where and how long users look on an interface. Large eye-tracking datasets of GUI visual search~\cite{vsgui10k} establish how factors such as visual clutter and target type affect \emph{search time}, and recent vision-language models predict scanpaths --- and, derived from them, search time --- given a screenshot together with a textual cue of the target to be found~\cite{seekui}. This work is closely related but differs in two respects central to the present proposal: it is \emph{target-specific} (it models the search for one named element) and it measures \emph{visual search time}, a component of, but not equal to, the time spent accomplishing a task on a screen.

\paragraph{The gap.} No existing approach predicts holistic, target-free per-screen Time-on-Task directly from a rendered screen, learned from heterogeneous real-world usage rather than constrained or laboratory settings. Classical models need manual operator coding and assume expert behaviour; learned performance models need structured UI inputs and constrained domains; visual-search models are target-cued and search-specific. This thesis addresses that gap.

\section{Research Question}
\begin{quote}
\itshape\color{accent}
To what extent can per-screen Time-on-Task be predicted from the visual and structural properties of a user interface screen alone --- independent of the user's goal, prior familiarity, and interaction history?
\end{quote}

\noindent
The question is deliberately stated at the level of the underlying construct rather than any particular model or dataset. Answering it well requires three things, which structure the work as \emph{objectives} rather than as separate questions:
\begin{enumerate}[label=(\roman*)]
  \item establishing a faithful procedure for deriving per-screen Time-on-Task labels from real interaction traces, with documented treatment of confounds such as system response latency;
  \item determining whether a learned predictor improves on simple, interpretable baselines, thereby quantifying the share of task time attributable to the screen itself; and
  \item testing whether a predictor trained on one source of timing data generalizes to human task-times measured independently, which separates a genuine relationship from an artifact of the training source.
\end{enumerate}

\section{Proposed Approach}
The approach is intentionally model-agnostic in its commitments; specific architectural and training choices are documented separately in an accompanying implementation plan and are not load-bearing for the research question.

\paragraph{Data and labels.} Training and evaluation use a large corpus of real-world web interaction traces~\cite{webchain} in which each step records a screen capture, its structural description, the action taken, and a timestamp. Per-screen Time-on-Task is derived as the interval between consecutive user actions within the same browser tab, after a verified-by-inspection account of what the timestamp marks. The label pipeline explicitly flags steps that involve navigation to a new page --- where the measured interval is partly system response latency rather than user effort --- so that the predictable, cognitive component can be analysed separately from confounded intervals. Applied to the corpus, the procedure yields on the order of $2\times10^5$ screen-level observations across roughly two hundred websites. Evaluation splits hold out entire registrable domains, so that no model is tested on a website it has seen during training, and every filtering and capping decision is logged in a public dataset card.

\paragraph{Models.} The central model reads a rendered screen and predicts its expected dwell time. To quantify how much the rendered pixels contribute, it is compared against an interpretable baseline built from countable structural features of the screen, each operationalizing a named construct of the classical performance-modelling canon: action type stands in for KLM's operator classes~\cite{klm1980}, interactive-element counts for the choice-difficulty term of the Hick--Hyman law~\cite{hick1952}, click-target size for the motor-difficulty term of Fitts' law~\cite{fitts1954}, and a navigation flag for KLM's system-response operator. Features enter in raw form and a gradient-boosted model learns each construct's functional shape, removing the manual operator coding that limits classical practice. Any margin the screen model achieves over this baseline is therefore, by construction, signal the classical constructs do not capture. A compact, openly available vision-language model is fine-tuned with parameter-efficient methods to keep the work reproducible on modest, shared computing resources; the emphasis is on the methodology and the comparison, not on identifying a best-performing model.

\paragraph{Validation against independent timing.} The trained predictor is finally evaluated, frozen and without any further tuning, on an independently collected dataset of cleanly measured human times: measured task times where available~\cite{tasksense}, with a public GUI visual-search timing dataset~\cite{vsgui10k} as the documented fallback --- noting explicitly that search time is one component of Time-on-Task, making the fallback a weaker but fully independent check. Agreement here --- in particular, agreement in the \emph{ranking} of screens by time cost --- is the strongest available evidence that the model captures a real property of interfaces rather than idiosyncrasies of its training source. This set is therefore quarantined: it is touched exactly once, a rule enforced mechanically in the experimental code rather than by convention alone.

\section{Evaluation}
Predictive accuracy is reported with standard regression error measures (in log-time and in seconds) on held-out interface domains, complemented by rank correlation between predicted and observed task times, since relative ordering of screens by cost is what a screening tool ultimately needs. The learned model is compared against constant-prediction floors and the interpretable structural baseline on identical evaluation rows. Results are broken down by whether a step involved navigation, isolating the component of task time that a screen can plausibly predict from the component driven by system latency. Calibration is examined to detect systematic over- or under-estimation. The independent-timing evaluation provides the headline test of generalization, reported as a rank correlation with a bootstrap confidence interval. Finally, predictions and measured human times are related to established computational models of perceptual load --- visual clutter metrics such as feature congestion and subband entropy~\cite{clutter,aim} --- asking how much of the learned model's estimate is reducible to quantified clutter and whether its residual still tracks human time; either outcome is informative about what a screen's predictable time cost consists of.

\section{Expected Contributions}
\begin{itemize}
  \item A validated, documented methodology for recovering per-screen Time-on-Task labels from real-world interaction logs, including the treatment of latency and exposure confounds.
  \item The first study, to our knowledge, of target-free per-screen Time-on-Task prediction directly from rendered interface screens on heterogeneous real-world data.
  \item An empirical characterisation of how much of human task time is recoverable from a screen alone, relative to interpretable structural baselines, and whether that signal generalizes to independently measured human timing.
\end{itemize}

\noindent
The work does not claim that a screen determines Time-on-Task: task time also depends on the user's goal, familiarity, and off-screen behaviour, which a single screen cannot observe. The contribution is to measure, rather than assume, the size of the recoverable component, and to do so with a method that scales far beyond manual modelling or user testing.

\section{Indicative Plan}
\begin{enumerate}[label=\arabic*.]
  \item Construct and document the label-extraction pipeline and dataset; report confound statistics.
  \item Implement the interpretable structural baseline as an early, self-contained result.
  \item Fine-tune and evaluate the screen-based predictor; compare against the baseline on held-out domains.
  \item Conduct the independent-timing generalization evaluation.
  \item Analysis and writing: confound breakdown, calibration, error analysis, and limitations.
\end{enumerate}

\begin{thebibliography}{99}
\small
\bibitem{hick1952} W.~E.~Hick. On the rate of gain of information. \emph{Quarterly Journal of Experimental Psychology}, 4(1):11--26, 1952.
\bibitem{fitts1954} P.~M.~Fitts. The information capacity of the human motor system in controlling the amplitude of movement. \emph{Journal of Experimental Psychology}, 47(6):381--391, 1954.
\bibitem{clutter} R.~Rosenholtz, Y.~Li, and L.~Nakano. Measuring visual clutter. \emph{Journal of Vision}, 7(2):17, 2007.
\bibitem{aim} A.~Oulasvirta et~al. Aalto Interface Metrics (AIM): A service and codebase for computational GUI evaluation. In \emph{Adjunct Proc.\ UIST}, 2018.
\bibitem{iso9241} International Organization for Standardization. \emph{ISO 9241-11:2018 --- Ergonomics of human-system interaction --- Part 11: Usability: Definitions and concepts.} ISO, 2018.
\bibitem{klm1980} S.~K.~Card, T.~P.~Moran, and A.~Newell. The keystroke-level model for user performance time with interactive systems. \emph{Communications of the ACM}, 23(7):396--410, 1980.
\bibitem{psychologyhci} S.~K.~Card, T.~P.~Moran, and A.~Newell. \emph{The Psychology of Human-Computer Interaction.} Lawrence Erlbaum Associates, 1983.
\bibitem{klmreview} A systematic review of modifications and validation methods for the extension of the Keystroke-Level Model. \emph{Advances in Human-Computer Interaction}, 2018.
\bibitem{deepui} P.~Duan et~al. Optimizing user interface layouts via gradient descent. In \emph{Proc.\ CHI}, 2020. arXiv:2002.10702.
\bibitem{vsgui10k} A.~Putkonen et~al. Understanding visual search in graphical user interfaces. \emph{International Journal of Human-Computer Studies}, 2025.
\bibitem{seekui} Y.~Guo, Y.~Jiang, L.~A.~Leiva, and A.~Oulasvirta. Predicting visual search behavior on graphical user interfaces with a reward-augmented vision-language model. In \emph{Proc.\ CHI}, 2026.
\bibitem{webchain} WebChain: A large-scale human-annotated dataset of real-world web interaction traces. 2026. arXiv:2603.05295.
\bibitem{tasksense} Yin et~al. Cognitive chain modelling and difficulty estimation for graphical user interface tasks. 2025. arXiv:2511.09309.
\end{thebibliography}

\end{document}